\documentclass[11pt]{article}

\usepackage{amsmath,amssymb,amsthm,mathrsfs}
\usepackage{geometry}
\geometry{margin=1in}
\usepackage{hyperref}
\newtheorem{theorem}{Theorem}
\begin{document}
\section{The Small-$a$ Limit of the Scattering Sector}

We analyze the behavior of the Morse scattering states as
\[
a \to 0,
\]
with $D_e$ fixed. Physically, this corresponds to flattening the
potential so that it approaches a constant background.

\subsection{Scaling of Parameters}

Recall the definitions
\[
\lambda = \frac{\sqrt{2mD_e}}{\hbar a},
\qquad
\kappa = \frac{k}{a},
\qquad
y = 2\lambda e^{-a(x-x_0)}.
\]

As $a\to 0$,
\[
\lambda \to \infty,
\qquad
\kappa \to \infty,
\qquad
\frac{\kappa}{\lambda}
=
\frac{\hbar k}{\sqrt{2mD_e}}
\quad \text{fixed}.
\]

Expand the potential:
\[
V(x)
=
D_e\left(1-e^{-a(x-x_0)}\right)^2.
\]

For small $a$,
\[
e^{-a(x-x_0)}
=
1-a(x-x_0)+\mathcal O(a^2),
\]
so
\[
V(x)
=
D_e a^2 (x-x_0)^2 + \mathcal O(a^3).
\]

Hence for fixed $x$,
\[
V(x)\to 0,
\]
while $V(+\infty)=D_e$ remains the asymptotic threshold.
Thus the Hamiltonian approaches
\[
H \to -\frac{\hbar^2}{2m}\frac{d^2}{dx^2} + D_e.
\]

\subsection{Limit of the Scattering Solutions}

The continuum wavefunctions are
\[
\psi_k(x)
=
\mathcal N(k)\,
y^{i\kappa}
e^{-y/2}
U\!\left(
\frac12-\lambda+i\kappa,
1+2i\kappa,
y
\right).
\]

We examine the asymptotics for large $\lambda,\kappa$.

Using
\[
y = 2\lambda e^{-a(x-x_0)}
=
2\lambda\left[1-a(x-x_0)+\mathcal O(a^2)\right],
\]
we write
\[
y = 2\lambda - 2\lambda a(x-x_0)+\mathcal O(a^2).
\]

Since
\[
\lambda a = \frac{\sqrt{2mD_e}}{\hbar},
\]
the combination
\[
2\lambda a(x-x_0)
=
\frac{2\sqrt{2mD_e}}{\hbar}(x-x_0)
\]
remains finite.

One uses the large-parameter asymptotics of the confluent
hypergeometric function $U(a,b,y)$
(see \cite{olver})
to obtain

\[
\psi_k(x)
\longrightarrow
\frac{1}{\sqrt{2\pi}}
e^{ikx},
\]
up to an irrelevant phase factor.

\subsection{Limit of the Reflection Amplitude}

Recall
\[
R(k)=
\frac{\Gamma(-2i\kappa)
\,\Gamma(\frac12-\lambda+i\kappa)}
{\Gamma(2i\kappa)
\,\Gamma(\frac12-\lambda-i\kappa)}.
\]

For large $\lambda,\kappa$,
using Stirling’s formula,

\[
\frac{\Gamma(\frac12-\lambda+i\kappa)}
{\Gamma(\frac12-\lambda-i\kappa)}
\sim
e^{2i\delta(k)},
\]
with phase shift

\[
\delta(k) \to 0
\quad (a\to 0).
\]

Hence
\[
R(k)\to 1.
\]

Since transmission is already absent,
this means scattering becomes trivial:
the potential becomes flat.

\subsection{Result}

\begin{theorem}
In the limit $a\to 0$ with $D_e$ fixed,
the Morse Hamiltonian converges (in strong resolvent sense)
to the free Hamiltonian shifted by $D_e$,
\[
H \to -\frac{\hbar^2}{2m}\frac{d^2}{dx^2}+D_e,
\]
and the scattering eigenfunctions converge to plane waves.
\end{theorem}

\subsection{Physical Interpretation}

As $a\to 0$:

\begin{itemize}
\item The exponential wall recedes to $x=-\infty$.
\item The potential becomes spatially flat.
\item Bound states disappear ($n_{\max}\to\infty$ but
      levels accumulate at $D_e$).
\item The S-matrix becomes trivial.
\end{itemize}

Thus the Morse potential interpolates between a confining
anharmonic well and a constant free background,
with $a$ controlling the geometric scale of the exponential wall.

\begin{thebibliography}{99}
\bibitem{olver}
F.~W.~J.~Olver et al.,
\textit{NIST Handbook of Mathematical Functions},
Cambridge Univ. Press (2010).
\end{thebibliography}


\section{The Hard-Wall Limit $a\to\infty$}

We now analyze the opposite regime,
\[
a \to \infty,
\]
with $D_e$ fixed.
This corresponds to compressing the exponential wall so that
the potential becomes infinitely steep at $x=x_0$.

\subsection{Pointwise Limit of the Potential}

Recall
\[
V(x)
=
D_e\left(1-e^{-a(x-x_0)}\right)^2.
\]

Consider separately the regions:

\paragraph{1. $x>x_0$.}
Then $e^{-a(x-x_0)}\to 0$ as $a\to\infty$,
so
\[
V(x)\to D_e.
\]

\paragraph{2. $x<x_0$.}
Then $e^{-a(x-x_0)}=e^{a(x_0-x)}\to\infty$,
so
\[
V(x)\to\infty.
\]

Thus, in the limit,

\[
V_\infty(x)
=
\begin{cases}
\infty, & x<x_0,\\[6pt]
D_e, & x>x_0.
\end{cases}
\]

This is a hard wall at $x=x_0$ together with a constant background.

\subsection{Limiting Hamiltonian}

The limiting operator is therefore

\[
H_\infty
=
-\frac{\hbar^2}{2m}\frac{d^2}{dx^2}
+
D_e,
\]
acting on $L^2(x_0,\infty)$
with Dirichlet boundary condition
\[
\psi(x_0)=0.
\]

\begin{theorem}
As $a\to\infty$, the Morse Hamiltonian converges
in the strong resolvent sense to
\[
H_\infty =
-\frac{\hbar^2}{2m}\frac{d^2}{dx^2}
+D_e
\quad \text{on } (x_0,\infty),
\]
with Dirichlet boundary condition at $x_0$.
\end{theorem}

\begin{proof}[Sketch]
The quadratic form of $H$ is
\[
Q_a[\psi]
=
\frac{\hbar^2}{2m}\int |\psi'|^2 dx
+
\int V(x)|\psi|^2 dx.
\]

As $a\to\infty$, $V(x)\to\infty$ for $x<x_0$,
forcing $\psi(x)=0$ in that region.
Hence the domain collapses to functions supported on $(x_0,\infty)$.
Standard results on monotone convergence of quadratic forms
imply strong resolvent convergence.
\end{proof}

\subsection{Continuous Spectrum}

The limiting spectrum is purely continuous:

\[
\sigma(H_\infty)
=
[D_e,\infty).
\]

The generalized eigenfunctions are

\[
\psi_k(x)
=
\sqrt{\frac{2}{\pi}}
\sin\big(k(x-x_0)\big),
\qquad
E=D_e+\frac{\hbar^2 k^2}{2m}.
\]

These are free half-line states satisfying Dirichlet boundary condition.

\subsection{Limit of the Reflection Amplitude}

Recall the Morse reflection amplitude:

\[
R(k)=
\frac{\Gamma(-2i\kappa)
\,\Gamma(\frac12-\lambda+i\kappa)}
{\Gamma(2i\kappa)
\,\Gamma(\frac12-\lambda-i\kappa)},
\qquad
\kappa=\frac{k}{a},
\qquad
\lambda=\frac{\sqrt{2mD_e}}{\hbar a}.
\]

As $a\to\infty$,

\[
\kappa\to 0,
\qquad
\lambda\to 0.
\]

Using small-argument expansion of Gamma functions,
one finds

\[
R(k)\to -1.
\]

Thus the phase shift tends to

\[
\delta(k)\to \frac{\pi}{2},
\]

which is precisely the phase shift for
a hard Dirichlet wall:

\[
e^{ikx} - e^{-ikx}
=
2i \sin(kx).
\]

\subsection{Discrete Spectrum in the Limit}

The bound-state energies are

\[
E_n
=
D_e
-
\frac{\hbar^2 a^2}{2m}
\left(\lambda-n-\frac12\right)^2.
\]

Since
\[
\lambda=\frac{\sqrt{2mD_e}}{\hbar a}\to 0,
\]
we have
\[
n_{\max}=\left\lfloor\lambda-\frac12\right\rfloor\to -1.
\]

Hence for sufficiently large $a$,
\[
\text{no bound states remain.}
\]

\subsection{Physical Interpretation}

As $a\to\infty$:

\begin{itemize}
\item The exponential wall sharpens into an impenetrable barrier.
\item The well structure disappears.
\item The discrete spectrum vanishes.
\item The system becomes a free particle on the half-line.
\item The reflection coefficient becomes $R=-1$.
\end{itemize}

Thus the Morse family interpolates between:

\[
\boxed{
\text{free line } (a\to 0)
\quad \longleftrightarrow \quad
\text{hard half-line } (a\to\infty)
}
\]

with intermediate values producing
a finite ladder of bound states and
nontrivial scattering phase.

\bigskip

This completes the description of the two extreme geometric limits
of the Morse Hamiltonian.

\section{Comparison of the Hard-Wall Limits:
Morse vs.\ P\"oschl--Teller}

We compare the limit $a\to\infty$ for the Morse potential
with the analogous large-parameter limit of the
hyperbolic P\"oschl--Teller potential.
Although both produce infinitely steep barriers,
their geometric realizations differ.

\subsection{1. Morse Hard-Wall Limit}

Recall
\[
V_M(x)=D_e\left(1-e^{-a(x-x_0)}\right)^2.
\]

As $a\to\infty$,

\[
V_M(x)\to
\begin{cases}
\infty, & x<x_0,\\[6pt]
D_e, & x>x_0.
\end{cases}
\]

The limiting Hamiltonian is

\[
H_\infty
=
-\frac{\hbar^2}{2m}\frac{d^2}{dx^2}
+D_e
\quad \text{on } (x_0,\infty)
\]

with Dirichlet boundary condition at $x_0$.
The spectrum is purely continuous:

\[
\sigma(H_\infty)=[D_e,\infty).
\]

The reflection coefficient tends to

\[
R(k)\to -1,
\]

corresponding to a perfectly reflecting wall.

\medskip

Thus the Morse family interpolates
\[
\mathbb R
\quad\longrightarrow\quad
\text{half-line }(x_0,\infty).
\]

\subsection{2. Hyperbolic P\"oschl--Teller Potential}

Consider the attractive hyperbolic P\"oschl--Teller potential

\[
V_{PT}(x)
=
-\frac{\hbar^2 a^2}{2m}
\frac{\lambda(\lambda-1)}{\cosh^2(ax)}.
\]

Its bound-state energies are

\[
E_n
=
-\frac{\hbar^2 a^2}{2m}
(\lambda-1-n)^2.
\]

The potential is symmetric and localized near $x=0$.

\subsection{Hardening the P\"oschl--Teller Well}

Two distinct limits may be considered:

\paragraph{(i) $a\to\infty$ with $\lambda$ fixed.}

Then
\[
\cosh^2(ax)\sim
\begin{cases}
1, & x=0,\\
\infty, & x\neq 0.
\end{cases}
\]

The potential collapses to a sharply localized well
at $x=0$,
approaching a delta-type interaction after suitable scaling.
The system remains defined on the full line.

\paragraph{(ii) $\lambda\to\infty$ with $a$ fixed.}

The well depth grows like $\lambda^2$.
After rescaling coordinates,
the system approaches an infinite square well on a finite interval.

Thus the P\"oschl--Teller family can produce
a two-sided confinement,
in contrast with Morse which produces a one-sided wall.

\subsection{3. Spectral Comparison}

\begin{center}
\begin{tabular}{lcc}
\hline
 & Morse ($a\to\infty$) & P\"oschl--Teller (large param.)\\
\hline
Domain & Half-line & Full line or finite interval \\
Discrete spectrum & disappears & may persist or diverge \\
Continuous spectrum & $[D_e,\infty)$ & $[0,\infty)$ \\
Reflection & $R=-1$ & symmetric scattering \\
Geometry & one-sided wall & symmetric well \\
\hline
\end{tabular}
\end{center}

\subsection{4. Algebraic Perspective}

The Morse system carries a truncated
$\mathfrak{su}(1,1)$ representation.
As $a\to\infty$,
the representation collapses and no discrete ladder remains.

The hyperbolic P\"oschl--Teller system also realizes
$\mathfrak{su}(1,1)$,
but in the deep-well limit the discrete representation
may become large rather than disappear.

The trigonometric P\"oschl--Teller system instead realizes
$\mathfrak{su}(2)$;
in its hard-wall limit it approaches an infinite square well,
retaining a purely discrete spectrum.

\subsection{5. Geometric Interpretation}

\begin{itemize}
\item Morse hard-wall limit:
      removes one half of configuration space.
\item Hyperbolic P\"oschl--Teller limit:
      produces symmetric short-range interaction.
\item Trigonometric P\"oschl--Teller limit:
      produces confinement on a compact interval.
\end{itemize}

Hence the Morse and P\"oschl--Teller hardening mechanisms
are geometrically distinct:

\[
\boxed{
\text{Morse: boundary formation}
\qquad
\text{P\"oschl--Teller: well deepening}
}
\]

\subsection{6. Confluent vs.\ Hypergeometric Origin}

Morse arises from a confluent hypergeometric equation,
while P\"oschl--Teller arises from the Gauss hypergeometric equation.

In the hard-wall limit:

\begin{itemize}
\item Morse $\rightarrow$ half-line free Laplacian.
\item Trigonometric P\"oschl--Teller $\rightarrow$ square well.
\item Hyperbolic P\"oschl--Teller $\rightarrow$ delta-type or deep symmetric well.
\end{itemize}

Thus the limiting operator retains memory of
the global geometry of the original system.

\bigskip

In conclusion, although both potentials admit
a steep-barrier limit,
the Morse family generates a geometric boundary,
whereas the P\"oschl--Teller family modifies the depth
or compactness of the confining region.

\end{document}
